\documentclass[12pt,a4paper]{article}
\usepackage{mathwizards-ingles}
\title{%
  \vspace{-1cm}
  {\Huge\bfseries\textcolor{mes1}{Mes 1 · Semanas 3--4}}\\[0.3cm]
  {\Large Rutinas Diarias, Present Simple y Adverbios de Frecuencia}\\[0.2cm]
  {\large\color{grissuave} Nivel A0+ $\rightarrow$ A1 · Transición al Reconocimiento de Patrones}\\[0.5cm]
  \rule{\textwidth}{1.5pt}
}
\author{}
\date{}

\begin{document}
\maketitle
\thispagestyle{empty}

\begin{cajanivel}
\textbf{Objetivo:} El estudiante comienza a reconocer patrones de rutina en el input. Todavía no se fuerza producción oral compleja, pero se empieza a tolerar respuestas de una o dos palabras. El foco está en \textbf{rutinas diarias}, el \textbf{Present Simple} como patrón emergente, y las \textbf{preguntas básicas}.
\end{cajanivel}

% ═══════════════════════════════════════════════════════════════════════════════
\section{Chunks: Rutinas Diarias (\textit{Daily Routines})}
% ═══════════════════════════════════════════════════════════════════════════════

\begin{cajachunk}[Morning Routine — Rutina Matutina]
\begin{tabularx}{\textwidth}{@{} X l @{}}
\eng{I wake up at 7 o'clock.} & \esp{Me despierto a las 7.} \\
\eng{I get up.} & \esp{Me levanto.} \\
\eng{I take a shower.} & \esp{Me ducho.} \\
\eng{I brush my teeth.} & \esp{Me lavo los dientes.} \\
\eng{I get dressed.} & \esp{Me visto.} \\
\eng{I have breakfast.} & \esp{Desayuno.} \\
\eng{I drink coffee.} & \esp{Bebo café.} \\
\eng{I go to work / school.} & \esp{Voy al trabajo / a la escuela.} \\
\end{tabularx}
\end{cajachunk}

\begin{cajachunk}[Afternoon / Evening Routine — Rutina Vespertina / Nocturna]
\begin{tabularx}{\textwidth}{@{} X l @{}}
\eng{I have lunch at noon.} & \esp{Almuerzo al mediodía.} \\
\eng{I come home.} & \esp{Llego a casa.} \\
\eng{I cook dinner.} & \esp{Cocino la cena.} \\
\eng{I watch TV.} & \esp{Veo televisión.} \\
\eng{I read a book.} & \esp{Leo un libro.} \\
\eng{I do my homework.} & \esp{Hago mi tarea.} \\
\eng{I go to bed at 10.} & \esp{Me acuesto a las 10.} \\
\eng{I fall asleep.} & \esp{Me duermo.} \\
\end{tabularx}
\end{cajachunk}

\begin{cajachunk}[Time Expressions — Expresiones de Tiempo]
\begin{multicols}{2}
\begin{tabularx}{\linewidth}{@{} X l @{}}
\eng{in the morning} & \esp{en la mañana} \\
\eng{in the afternoon} & \esp{en la tarde} \\
\eng{in the evening} & \esp{en la noche} \\
\eng{at night} & \esp{de noche} \\
\end{tabularx}

\columnbreak

\begin{tabularx}{\linewidth}{@{} X l @{}}
\eng{every day} & \esp{todos los días} \\
\eng{every morning} & \esp{todas las mañanas} \\
\eng{at \underline{\hspace{1cm}} o'clock} & \esp{a las \underline{\hspace{1cm}}} \\
\eng{on Monday(s)} & \esp{los lunes} \\
\end{tabularx}
\end{multicols}
\end{cajachunk}

% ═══════════════════════════════════════════════════════════════════════════════
\section{Gramática Emergente: Present Simple}
% ═══════════════════════════════════════════════════════════════════════════════

\begin{cajagrammar}[Present Simple — Afirmativo]
Se usa para hablar de \textbf{rutinas, hábitos y hechos generales}. Es el primer tiempo verbal que el estudiante reconocerá en el input.

\medskip
\begin{tabularx}{\textwidth}{@{} l l X @{}}
\toprule
\textbf{Sujeto} & \textbf{Verbo} & \textbf{Ejemplo} \\
\midrule
I / You / We / They & verbo base & \textit{I \textbf{wake up} at 7.} \\
He / She / It & verbo + \textbf{-s} & \textit{She \textbf{wakes} up at 7.} \\
\bottomrule
\end{tabularx}

\medskip
\textbf{Regla de la ``s'':} En tercera persona singular (he, she, it), se agrega \textbf{-s} al verbo.

\begin{cajaejemplo}[Patrones de la tercera persona]
\begin{tabularx}{\textwidth}{@{} l l l @{}}
\textbf{Regla} & \textbf{Base} & \textbf{He/She/It} \\
\midrule
General: +s & \textit{eat} & \textit{eat\textbf{s}} \\
Termina en -s, -sh, -ch, -x, -o: +es & \textit{wash} & \textit{wash\textbf{es}} \\
Termina en consonante + y: -y $\rightarrow$ -ies & \textit{study} & \textit{stud\textbf{ies}} \\
Irregular & \textit{have} & \textit{\textbf{has}} \\
\end{tabularx}
\end{cajaejemplo}
\end{cajagrammar}

\begin{cajagrammar}[Present Simple — Negativo e Interrogativo]
\begin{tabularx}{\textwidth}{@{} l X @{}}
\toprule
\textbf{Forma} & \textbf{Estructura} \\
\midrule
Negativo (I/you/we/they) & \textit{I \textbf{don't} (do not) eat breakfast.} \\
Negativo (he/she/it) & \textit{She \textbf{doesn't} (does not) eat breakfast.} \\
Pregunta (I/you/we/they) & \textit{\textbf{Do} you eat breakfast?} \\
Pregunta (he/she/it) & \textit{\textbf{Does} she eat breakfast?} \\
\bottomrule
\end{tabularx}

\medskip
\begin{cajaejemplo}[Nota Importante]
Después de \textbf{does / doesn't}, el verbo vuelve a su forma base (sin la ``s''):\\
$\checkmark$ \textit{She \textbf{doesn't} \textbf{eat}} \quad $\times$ \textit{She doesn't eat\textbf{s}}
\end{cajaejemplo}
\end{cajagrammar}

% ═══════════════════════════════════════════════════════════════════════════════
\section{Chunks: Adverbios de Frecuencia}
% ═══════════════════════════════════════════════════════════════════════════════

\begin{cajachunk}[Frequency Adverbs — Posición y Significado]

\begin{center}
\begin{tabularx}{0.8\textwidth}{@{} l c X @{}}
\toprule
\textbf{Adverbio} & \textbf{Frecuencia} & \textbf{Ejemplo} \\
\midrule
\eng{always} & 100\% & \textit{I \textbf{always} drink coffee.} \\
\eng{usually} & $\sim$80\% & \textit{She \textbf{usually} walks to school.} \\
\eng{often} & $\sim$60\% & \textit{We \textbf{often} watch movies.} \\
\eng{sometimes} & $\sim$40\% & \textit{He \textbf{sometimes} reads at night.} \\
\eng{rarely} / \eng{seldom} & $\sim$10\% & \textit{They \textbf{rarely} eat out.} \\
\eng{never} & 0\% & \textit{I \textbf{never} wake up early.} \\
\bottomrule
\end{tabularx}
\end{center}

\medskip
\textbf{Posición:} Los adverbios de frecuencia van \textbf{antes del verbo principal} pero \textbf{después del verbo \textit{to be}}:
\begin{itemize}[nosep, leftmargin=1cm]
  \item $\checkmark$ \textit{I \textbf{always} eat breakfast.} (antes de ``eat'')
  \item $\checkmark$ \textit{She \textbf{is} \textbf{always} happy.} (después de ``is'')
\end{itemize}
\end{cajachunk}

% ═══════════════════════════════════════════════════════════════════════════════
\section{Chunks: Preguntas Básicas (\textit{Question Words})}
% ═══════════════════════════════════════════════════════════════════════════════

\begin{cajachunk}[WH- Questions — Las Preguntas Fundamentales]
\begin{tabularx}{\textwidth}{@{} l l X @{}}
\toprule
\textbf{Palabra} & \textbf{Significado} & \textbf{Ejemplo} \\
\midrule
\eng{What} & Qué & \textit{\textbf{What} do you eat for breakfast?} \\
\eng{Where} & Dónde & \textit{\textbf{Where} do you live?} \\
\eng{When} & Cuándo & \textit{\textbf{When} do you wake up?} \\
\eng{Who} & Quién & \textit{\textbf{Who} is your teacher?} \\
\eng{Why} & Por qué & \textit{\textbf{Why} do you study English?} \\
\eng{How} & Cómo & \textit{\textbf{How} do you go to school?} \\
\eng{How many} & Cuántos & \textit{\textbf{How many} books do you have?} \\
\eng{What time} & A qué hora & \textit{\textbf{What time} do you wake up?} \\
\bottomrule
\end{tabularx}
\end{cajachunk}

% ═══════════════════════════════════════════════════════════════════════════════
\section{Chunks: Preposiciones de Lugar}
% ═══════════════════════════════════════════════════════════════════════════════

\begin{cajachunk}[Prepositions of Place — Preposiciones de Lugar]
\begin{multicols}{2}
\begin{tabularx}{\linewidth}{@{} l l @{}}
\eng{in} & \esp{en, dentro de} \\
\eng{on} & \esp{sobre, encima de} \\
\eng{under} & \esp{debajo de} \\
\eng{next to} & \esp{al lado de} \\
\eng{between} & \esp{entre} \\
\eng{behind} & \esp{detrás de} \\
\eng{in front of} & \esp{delante de} \\
\eng{near} & \esp{cerca de} \\
\end{tabularx}

\columnbreak

\begin{cajaejemplo}[Ejemplos]
\begin{itemize}[nosep, leftmargin=0.5cm]
  \item \textit{The cat is \textbf{under} the table.}
  \item \textit{My phone is \textbf{on} the desk.}
  \item \textit{The park is \textbf{next to} the school.}
  \item \textit{The dog is \textbf{behind} the door.}
  \item \textit{She lives \textbf{near} the hospital.}
\end{itemize}
\end{cajaejemplo}
\end{multicols}
\end{cajachunk}

% ═══════════════════════════════════════════════════════════════════════════════
\section{Vocabulario Temático: La Casa y Objetos Cotidianos}
% ═══════════════════════════════════════════════════════════════════════════════

\begin{cajachunk}[Around the House — Vocabulario del Hogar]
\begin{multicols}{3}
\begin{itemize}[nosep, leftmargin=0.5cm]
  \item \eng{house} \esp{casa}
  \item \eng{room} \esp{habitación}
  \item \eng{kitchen} \esp{cocina}
  \item \eng{bathroom} \esp{baño}
  \item \eng{bedroom} \esp{dormitorio}
  \item \eng{living room} \esp{sala}
  \item \eng{door} \esp{puerta}
  \item \eng{window} \esp{ventana}
  \item \eng{table} \esp{mesa}
  \item \eng{chair} \esp{silla}
  \item \eng{bed} \esp{cama}
  \item \eng{desk} \esp{escritorio}
  \item \eng{phone} \esp{teléfono}
  \item \eng{computer} \esp{computadora}
  \item \eng{book} \esp{libro}
  \item \eng{pen} \esp{bolígrafo}
  \item \eng{cup} \esp{taza}
  \item \eng{plate} \esp{plato}
\end{itemize}
\end{multicols}
\end{cajachunk}

% ═══════════════════════════════════════════════════════════════════════════════
\section{Oraciones Listas para Minar}
% ═══════════════════════════════════════════════════════════════════════════════

\begin{cajamineria}[Oraciones \textit{i+1} — Semanas 3--4]

\begin{enumerate}[leftmargin=1.2cm]
  \item \textit{I \underline{usually} wake up at seven o'clock.}\\
    {\small\textbf{Objetivo:} el adverbio ``usually'' (usualmente).}

  \item \textit{She \underline{brushes} her teeth every morning.}\\
    {\small\textbf{Objetivo:} la regla de la ``-es'' en tercera persona (brush $\rightarrow$ brushes).}

  \item \textit{Do you \underline{have breakfast} every day?}\\
    {\small\textbf{Objetivo:} la colocación ``have breakfast'' (desayunar).}

  \item \textit{My keys are \underline{on} the table.}\\
    {\small\textbf{Objetivo:} la preposición ``on'' (sobre).}

  \item \textit{He \underline{never} drinks coffee. He prefers tea.}\\
    {\small\textbf{Objetivo:} el adverbio ``never'' y su posición antes del verbo.}

  \item \textit{\underline{What time} does she go to bed?}\\
    {\small\textbf{Objetivo:} la expresión interrogativa ``what time'' (¿a qué hora?).}

  \item \textit{I don't eat meat. I \underline{always} eat vegetables.}\\
    {\small\textbf{Objetivo:} la estructura negativa ``don't'' + contraste con ``always''.}

  \item \textit{The cat is \underline{under} the chair.}\\
    {\small\textbf{Objetivo:} la preposición ``under'' (debajo de).}

  \item \textit{We \underline{sometimes} go to the park on Sundays.}\\
    {\small\textbf{Objetivo:} ``sometimes'' y su posición en la oración.}

  \item \textit{Where does your \underline{brother} live?}\\
    {\small\textbf{Objetivo:} ``brother'' (hermano) + reconocer la estructura con ``does''.}

  \item \textit{She \underline{studies} English every afternoon.}\\
    {\small\textbf{Objetivo:} la regla consonante + y $\rightarrow$ -ies (study $\rightarrow$ studies).}

  \item \textit{I \underline{come home} at five o'clock.}\\
    {\small\textbf{Objetivo:} la expresión ``come home'' (llegar a casa) sin preposición.}
\end{enumerate}
\end{cajamineria}

\vfill
\begin{center}
\rule{0.5\textwidth}{0.5pt}\\[0.3cm]
{\small\color{grissuave} Curso de Inmersión en Inglés · Mes 1 · Semanas 3--4 · Nivel A0+ $\rightarrow$ A1}
\end{center}

\end{document}
